\documentclass[a4paper]{article}
\usepackage[pdftex]{graphicx}
\usepackage[utf8]{inputenc}
\usepackage{enumerate}
\usepackage{amssymb}
\usepackage{icomma}
\usepackage{siunitx}
\sisetup{locale=DE}
\usepackage{tikz}
\usepackage{href-ul}
\hypersetup{
	colorlinks=true,
	linkcolor=blue,
	urlcolor=blue}
\usepackage{geometry}
\geometry{a4paper, top=15mm, left=15mm, right=15mm, bottom=15mm,
	headsep=10mm, footskip=12mm}

\begin{document}
	\begin{tikzpicture}(10,3)
		\draw[ultra thick](2,0) --(10,0) -- (10,1.3) --(2,1.3) -- (2,0);
		\draw[fill=black](2,0)-- (0,.6) -- (2,.6) -- (2,0);
		\node at (6,.6) {\large solution sheet from \href{http://www.okuyakl.com}{www.okuyakl.com}};
	\end{tikzpicture}
	\vspace{0.5 cm}
	
	\noindent
	\begin{minipage}{0.5\textwidth}
		\noindent{\bf Task 1. a)}\\
		\includegraphics[width=7 cm]{tangent052}
	\end{minipage}
	\hfill
	\begin{minipage}{0.5\textwidth}
		\noindent{\bf Task 1. b)}\\
		The local line is the Thales circle around $M$ through $P$ and $A$.
	\end{minipage}
	
	\noindent{\bf Task 2.}\\
	These are two parallels that form a pair and are at a distance of $\SI{8}{\centi\meter}$ from the central parallel $g$:
	$$M=\{ P| d(P;g)=\SI{8}{\centi\meter}\}$$
	
	\noindent {\bf task 3. a)}\\
	$$A=\{ P|d(P;M_1)>r_1 \land P| d(P;M_2)<r_2 \} $$
	
	\noindent {\bf task 3. b)}\\
	$$B=\{ P|d(P;g)\ge d(P;h)\} $$
	
	\noindent {\bf task 3. c)}\\
	The angle between $w_1$ and $w_2$ is $\alpha + \beta$.
	For the angles $\alpha$ and $\beta$ the following applies:
	$$
	\renewcommand{\arraystretch}{1.5}
	\begin{array}{rcll}
		2\cdot \alpha + 2 \cdot \beta &=& 180^\circ & |:2 \\
		\alpha + \beta &=& 90^\circ
		
	\end{array}
	$$
	So $w_1$ and $w_2$ form a right angle.
	
	\noindent {\bf Task 4. a)}\\
	$M_1$ are two locus lines, namely the bisectors of $g_1$ and $g_2$:
	$$M_1=\{P|d(P;g_1)=d(P;g_2)\}$$
	
	\noindent {\bf task 4. b)}\\
	$M_2$ is a local line, the central parallel between $h_1$ and $h_2$:
	$$M_2=\{P|d(P;h_1)=d(P;h_2)\}$$
	
	\noindent {\bf task 4. b)}\\
	$K$ is a spatial region, namely the interior of the circle $k(M;r)$ around $M$ with radius $r$:
	$$K=\{P|d(P;M)\le r\}$$
	
	\noindent {\bf task 5.)}\\
	
	\begin{minipage}{0.5\textwidth}
		\includegraphics[width=5 cm]{papa160}
	\end{minipage}
	\hfill
	\begin{minipage}{0.5\textwidth}
		P lies on a pair of parallel lines at a distance of $\SI{1.5}{\centi\meter}$ from $g$
	\end{minipage}
	
	\noindent {\bf task 6.a)}\\
	The set of points is a location line because the distance always has a certain value.
	It describes a circular line around $M$ with radius $\SI{2.8}{\centi\meter}$.
	
	\noindent {\bf Task 6.b)}\\
	This point set is a spatial region because the distance is greater than or equal to $r_2$. She
	describes the area outside a circle around $M$ with radius $r_2$.
	
	\noindent {\bf task 7.)}\\
	\begin{minipage}{0.6\textwidth}
		\includegraphics[width=6 cm]{pum160}
	\end{minipage}
	\hfill
	\begin{minipage}{0.4\textwidth}
		{\bf a)} P lies within a Thales circle between $A$ and $C$\\
		{\bf b)} P' lies on and to the right of the perpendicular bisector of $A$ and $B$.
	\end{minipage}
	
	\noindent {\bf task 8.)}\\
	$$G=\{ P| \SI{2}{\centi \meter} \le d(P;M) < \SI{4}{\centi \meter} \}$$
	
	\begin{center}
		\includegraphics[width=7 cm]{../../viecher/eendcomic.pdf}
		
		Here you can go back to the \href{https://www.okuyakl.de/math/m8ortL052/ae052.pdf}{task sheet}
	\end{center}
	
\end{document}